Data plays a critical role in machine learning. Every machine learning
model is trained and evaluated using data, quite often in the
form \edit{of} static dataset\edit{s}. The characteristics of these
datasets fundamentally influence a model's behavior: \edit{a} model is
unlikely to perform well in the wild if its deployment context does
not match its training or evaluation datasets, or if these datasets
reflect unwanted
\edit{societal} biases. Mismatches like this can have especially severe
consequences when machine learning \edit{models are} used in
high-stakes domains\edit{,} such as criminal
justice~\cite{facerec_law_us, facerec_law_elsewhere, compass},
hiring~\cite{hiring_ml_us}, critical infrastructure~\cite{water,
grid}, \edit{and} finance~\cite{newinvestor}. \edit{E}ven in other
domains, mismatches may lead to loss of revenue or public relations
setbacks. Of particular concern are recent examples showing that
machine learning models can reproduce or amplify unwanted societal
biases reflected in training \edit{datasets} \cite{gendershades,
amazon, word2vec}. For these and other reasons, the World Economic
Forum suggests that all entities should document the provenance,
creation, and use of machine learning datasets in order to avoid
discriminatory outcomes~\cite{wef}.

Although data provenance has been studied extensively in the databases
community~\cite{provenance,datahub}, it is rarely discussed in the
machine learning community\edit{. D}ocumenting the creation and use of
datasets has received even less attention. Despite the importance of
data to machine learning, there is \edit{currently no} standardized
process for documenting machine learning datasets.

To address this gap, we propose \emph{datasheets for datasets}. In the
electronics industry, every component, no matter how simple or
complex, is accompanied with a datasheet describing its operating
characteristics, test results, recommended usage, and other
information. By analogy, we propose that every dataset be accompanied
with a datasheet that documents its motivation, composition,
collection process, recommended uses, and so on. Datasheets for
datasets have the potential to increase transparency and
accountability within the machine learning community, mitigate
unwanted \edit{societal} biases in machine learning \edit{models},
facilitate greater reproducibility of machine learning results, and
help researchers and practitioners \edit{to} select more appropriate datasets
for their chosen tasks.

After outlining our objectives below, we describe the process by which
we developed \edit{datasheets for datasets}. \edit{We then provide a set of
questions designed to elicit the information that a datasheet for a
dataset might contain, as well as a} workflow for dataset creators to
use when answering these questions. We conclude with a summary of the
impact \edit{to date} of datasheets for datasets and a discussion of
implementation challenges and avenues for future work.
